\section*{Задача 2. Скупой рыцарь}
Скупой рыцарь хранит в подвале множество сундуков, занумерованных начиная с единицы. 
В каждом сундуке лежат монеты. О содержимом каждого сундука есть опись. 
В описи используются строчные латинские буквы:
\begin{itemize}
  \item \texttt{a} --- $1$ копейка, \texttt{b} --- $5$ копеек, \texttt{c} --- $10$ копеек,
  \item \texttt{d} --- $50$ копеек, \texttt{e} --- $1$ рубль, \texttt{f} --- $2$ рубля,
  \item \texttt{g} --- $5$ рублей, \texttt{h} --- $10$ рублей, \texttt{i} --- $25$ рублей.
\end{itemize}
Буква входит в опись столько раз, сколько монет соответствующего номинала лежит в сундуке. 
Буквы могут идти в произвольном порядке. Опись пустого сундука --- пустая строка.

Рыцарь хочет найти два сундука $K$ и $L$ ($K < L$), такие что модуль разности денежных сумм 
$\left| S_K - S_L \right|$ максимален. Если таких пар несколько, выбирается та, у которой 
сумма номеров $K+L$ максимальна.

\textbf{Формат ввода:}\\
В первой строке число $N$ ($2 \le N \le 1000$) --- количество сундуков. 
В следующих $N$ строках содержатся описи сундуков (длины от $0$ до $1000$ символов).

\textbf{Формат вывода:}\\
В первой строке выведите $K$, во второй --- $L$.

\begin{examplebox}
  \begin{minipage}[t]{0.48\linewidth}
    \textbf{Ввод:}\\
    \begin{verbatim}
2
abbahihihi
cdffge
    \end{verbatim}
  \end{minipage}%
  \hfill%
  \begin{minipage}[t]{0.48\linewidth}
    \textbf{Вывод:}\\
    \begin{verbatim}
1
2
    \end{verbatim}
  \end{minipage}
\end{examplebox}